\section*{Bab \ref{ch:g3:animals-through-seasons} --- Hewan Sepanjang Musim}

\begin{solution}{exo:g3:animals-through-seasons:1}
Dingin, yang membuat menjaga diri tetap hangat menghabiskan lebih
banyak \omterm{prop:g3:balanced-diet:jobs}{energi} --- tepat pada saat makanan, sumber \omterm{prop:g3:balanced-diet:jobs}{energi} itu, menjadi
langka.
\end{solution}

\begin{solution}{exo:g3:animals-through-seasons:2}
\omterm{def:g3:animals-through-seasons:migration}{Bermigrasi}: \omterm{prop:g3:sorting-living-things:groups}{burung} layang-layang, bangau (atau \omterm{prop:g3:sorting-living-things:groups}{burung} jenjang).
\omterm{def:g3:animals-through-seasons:hibernation}{Berhibernasi}: landak, marmot gunung (atau tikus \omterm{prop:g1:healthy-habits:needs}{tidur}, kelelawar).
\omterm{def:g3:animals-through-seasons:active}{Tetap aktif}: rubah, tupai (atau \omterm{prop:g3:sorting-living-things:groups}{burung} robin, rusa).
\end{solution}

\begin{solution}{exo:g3:animals-through-seasons:3}
Tubuhnya mendingin, jantung dan napasnya melambat sampai hampir tidak
ada, dan ia \omterm{prop:g1:healthy-habits:needs}{tidur} nyenyak di tempat berlindungnya. Ia hidup dari lemak
yang ditimbunnya pada musim gugur.
\end{solution}

\begin{solution}{exo:g3:animals-through-seasons:4}
\omterm{prop:g3:sorting-living-things:groups}{Burung} layang-layang memakan \omterm{prop:g3:sorting-living-things:groups}{serangga} terbang yang ditangkapnya sambil
melayang --- dan pada musim dingin di sini \omterm{prop:g3:sorting-living-things:groups}{serangga} itu sama sekali
tidak ada. Karena persediaan makanannya lenyap sama sekali, ia harus
pergi ke tempat \omterm{prop:g3:sorting-living-things:groups}{serangga} masih terbang: negeri yang lebih hangat.
\end{solution}

\begin{solution}{exo:g3:animals-through-seasons:5}
\omterm{def:g3:animals-through-seasons:active}{Mantel musim dingin} yang lebih tebal (rubah, rusa); \omterm{def:g3:animals-through-seasons:active}{cadangan makanan}
yang dikumpulkan pada musim gugur (kacang yang dikubur tupai); daftar
makanan yang berubah (\omterm{prop:g3:sorting-living-things:groups}{burung} hitam yang beralih dari cacing ke \omterm{prop:g3:flowers-fruits-seeds:transform}{buah}
beri); dan pencarian yang lebih keras (\omterm{prop:g3:sorting-living-things:groups}{burung} robin yang berpatroli
mencari \omterm{prop:g3:flowers-fruits-seeds:transform}{buah} beri terakhir).
\end{solution}

\begin{solution}{exo:g3:animals-through-seasons:6}
Si \omterm{def:g3:animals-through-seasons:migration}{pemigrasi} menanggung risiko perjalanan yang melelahkan; si
penghibernasi harus mengumpulkan seluruh lemak musim dinginnya pada
musim gugur atau tidak terbangun pada musim semi; \omterm{def:g1:animals-around-us:animal}{hewan} yang aktif
bertaruh pada penemuan makanan sepanjang bulan-bulan yang paling
lapar.
\end{solution}

\begin{solution}{exo:g3:animals-through-seasons:7}
Hibernasi --- makan besar-besaran membangun lemak musim dinginnya,
lubang berlapis \omterm{def:g1:plants-around-us:parts}{daun} itu sarang musim dinginnya. Langkah 4: pemeriksaan
musim gugur (menggemuk dan menggali sarang menunjuk pada hibernasi).
\end{solution}

\begin{solution}{exo:g3:animals-through-seasons:8}
Jawaban terbuka. Tidak ada \omterm{prop:g3:sorting-living-things:groups}{burung} layang-layang pada musim dingin ---
mereka berada di Afrika. \omterm{prop:g3:sorting-living-things:groups}{Burung} yang sungguh-sungguh terlihat (robin,
gelatik, \omterm{prop:g3:sorting-living-things:groups}{burung} hitam, pipit, gagak) adalah \omterm{prop:g3:sorting-living-things:groups}{burung} bersiasat aktif.
\end{solution}

\begin{solution}{exo:g3:animals-through-seasons:9}
(a) Makanan tupai masih dapat ditemukan pada musim dingin --- terutama
yang dikuburnya pada musim gugur --- jadi: aktif, hidup dari
cadangannya. (b) Katak dan \omterm{prop:g3:sorting-living-things:groups}{serangga} \omterm{prop:g3:sorting-living-things:groups}{burung} jenjang lenyap dari padang
rumput basah yang membeku --- makanannya hilang, jadi: \omterm{def:g3:animals-through-seasons:migration}{migrasi} (\omterm{prop:g3:sorting-living-things:groups}{burung}
jenjang memang terbang ke selatan dalam kawanan besar berbentuk V).
\end{solution}

\begin{solution}{exo:g3:animals-through-seasons:10}
Seorang penghibernasi membayar musim dinginnya di muka, dengan lemak
musim gugur; landak yang kurus tidak punya apa pun untuk dibakar
sepanjang bulan-bulan \omterm{prop:g1:healthy-habits:needs}{tidur} dinginnya dan tidak akan terbangun. Para
penyelamat menampung landak musim gugur yang terlalu kurus lalu
memberinya makan di dalam ruangan sampai beratnya cukup untuk
\omterm{def:g3:animals-through-seasons:hibernation}{berhibernasi} dengan aman --- atau menjaganya tetap hangat dan
kenyang sepanjang musim dingin.
\end{solution}
